\expandafter\def\csname HistoricalNoteText491\endcsname{\emph{Ueber die Berechnung der Cometenbahnen.} \emph{Von} \textsc{Franz Güssmann}. \emph{Wien, bei Joh. Thom. Edlen von Trattnern}, 1803. 4to. 43 Seiten mit einer Kupferplatte. Eigentlich hätte eine so armselige, mit Kuttenwitz ausgeschmückte Schmähschrift wohl überhaupt keine Erwähnung verdient. Es verdient aber Beachtung, dass man einen Menschen in seiner wahren Gestalt und Blöße zeigen sollte, der nur umherschleicht, um verdienten Männern bei nicht hinreichend unterrichteten Personen zu schaden. Dieser neidische und hämische Jesuit hat von jeher nichts als Schmähschriften \emph{à la Valencia} geschrieben. Ein kleines Beispiel oben zeigt, wie gewissenhaft und einsichtsvoll dieser heilige und gelehrte Mann dabei vorgeht. Wir könnten noch mehrere solcher Beispiele anführen, wenn es die Mühe lohnte.}
\clearpage
\fancyhead[L]{\small GÜSSMANN. KOMETENBAHNEN.}
\hypertarget{native-491-lower}{}\label{native-491-lower}
\begin{center}
\small Monatliche Correspondenz. Band XV. S. 452..460. Mai 1807.\\[3pt]
\rule{.72\linewidth}{.3pt}
\end{center}
\begin{center}Auszug aus einem Brief.\end{center}
\begin{flushright}--- --- --- 10. März 1807.\end{flushright}

--- --- --- Ich habe die Ehre, Ihnen die gewünschte Broschüre\HistoricalNote{491} von \textsc{Güssmann} zu übersen%
\SourceLeaf{492}%
den. Eine ausführliche und detaillierte Antwort ist sie wohl nicht wert. Eine solche Antwort könnte nur interessant werden, wenn man die Methoden von \textsc{Olbers} und \textsc{Boscovich} miteinander vergliche und zeigte, welche Vorteile die erstere gegenüber der letzteren durch ihre größere Bequemlichkeit hat. Man kann von diesen Vorteilen durchaus überzeugt sein, ohne \textsc{Boscovich}s Ruhm zu schmälern. Von diesem Ruhm kommt seinem Ordensbruder natürlich nichts zugute. Die Würdigung der Broschüre des Letzteren, die sich keineswegs dafür eignet, in diese Untersuchung einbezogen zu werden, bleibt daher wohl besser davon getrennt.

Auch nur mit einem Wort auf die ungezogenen, armseligen und völlig sachfremden Einfälle einzugehen, wäre sowohl Ihnen als auch Dr. \textsc{Olbers} gegenüber höchst unangebracht. Ich begnüge mich damit, die Hauptpunkte herauszuheben, mit denen er Unkundigen etwas gegen die Methode von Dr. \textsc{Olbers} gesagt zu haben scheinen könnte.

\textsc{Güssmann}s wichtigster Angriff richtet sich gegen die Formel von Dr. \textsc{Olbers}, mit der das Verhältnis der kurtierten Abstände in der ersten und dritten Beobachtung bestimmt wird. \textsc{Güssmann} gibt dann \emph{auf seine Art} eine Herleitung, in der er eine Menge Ungereimtheiten und Fehler findet. Dr. \textsc{Olbers} wird eine solche Herleitung jedoch nicht als seine eigene anerkennen, und wer dessen klare Darstellung mit \textsc{Güssmann}s sonderbarer Begriffsverwirrung vergleicht, wird zwischen beiden keinerlei Übereinstimmung finden. Ich weiß nicht, ob ich diese durchgängige Entstellung von \textsc{Olbers}’ Darstellung für \emph{absichtlich} halten oder glauben soll, dass \textsc{Güssmann} sie selbst nicht verstanden hat. Schon in den §§ 1 und 2 und später in § 21 wird \textsc{Olbers} die Annahme \emph{angedichtet}, die Bewegung des Kometen werde als geradlinig angesehen. Zwischen der einzigen Annahme von Dr. \textsc{Olbers}, dass die Sehnen der zwischen der ersten und der letzten Beobachtung durchlaufenen Bögen der Kometen- und Erdbahn von den mittleren \emph{Radiusvektoren} im Verhältnis der Zeiten geteilt werden, und der Annahme, die \textsc{Güssmann} ihm \emph{aufdrängen} will, dass sich der Komet und die Erde auf den Sehnen selbst bewegen, besteht ein so großer Unterschied, dass Gedankenlosigkeit oder — beispiellose Unverschämtheit dazugehört, ihn unwissentlich oder wissentlich zu übersehen. \textsc{Olbers}’ Herleitung des oben genannten Verhältnisses ($M$) ist so klar, wie es jeder Leser, für den seine Schrift bestimmt ist, erwarten kann. Es seien $k$ und $i$ die äußeren geozentrischen Positionen des Kometen auf der Himmelskugel. $h$ sei die Position, an der der Komet von der Erde aus erscheinen würde, wenn beide auf den Sehnen zwischen ihren äußeren Positionen an den Punkten stünden, an denen die mittleren \emph{Radiusvektoren} diese Sehnen schneiden. $h^*$ sei die wirkliche geozentrische Position bei der mittleren Beobachtung auf der Himmelskugel, $g$ die mittlere Position der Sonne. Dann liegen einerseits $g$, $h$ und $h^*$ auf einem Großkreis und andererseits auch $k$, $h$ und $i$. Ersteres gilt, weil die gedachte Sichtlinie offensichtlich in derselben Ebene wie die Sonne und die mittleren Positionen der Erde und des Kometen liegt. Das Zweite ergibt sich leicht — und \textsc{Olbers} hat es zu Recht seinen Lesern überlassen, dies herzuleiten — aus der Annahme, dass die Abschnitte, in die die Sehnen der Erdbahn und der Planetenbahn von den mittleren \emph{Radiusvektoren} geteilt werden, einander proportional sind. Darauf gründet \textsc{Olbers} nun ein einfaches Verfahren zur Bestimmung der Lage von $h$ auf der Himmelskugel. Damit wird dann das Verhältnis der kurtierten Abstände in den äußeren Beobachtungen, $M$, bestimmt. Schließlich zeigt \textsc{Olbers}, wie man auch auf die vorläufige Bestimmung von
\SourceLeaf{493}
$h$ ganz verzichten und $M$ unmittelbar aus den Daten der Aufgabe berechnen kann. Dieser letzte Schritt hätte analytisch mit den bereits gefundenen Formeln ausgeführt werden können. \textsc{Olbers} hielt es jedoch für anschaulicher, ihn mithilfe einer Projektion herzuleiten. Man vergleiche nun diese klare Darstellung mit \textsc{Güssmann}s abenteuerlicher Verwirrung. In § 6 und in Fig. II bezeichnet er wie ich die äußeren Positionen mit $k$ und $i$, mit $h$ aber zunächst die \emph{wahre Position}. Er behauptet, \textsc{Olbers} setze voraus, $h$ liege mit $k$ und $i$ auf demselben Großkreis. Das ist abermals eine unverschämte Unwahrheit, denn das hat \textsc{Olbers} nirgends gesagt. Gleich danach verwendet er $h$ jedoch für die oben so bezeichnete gedachte Position. Diese Verwechslung bringt ihn nun völlig durcheinander. Er meint, nach \textsc{Olbers}’ Verfahren müsse sich für $c''$, das heißt für die Länge des Punktes $h$, immer wieder die Länge des Punktes $h^*$, die mit $a''$ bezeichnet ist, ergeben. Zu diesem sonderbaren Urteil haben ihn nur die \emph{willkürliche} Verwechslung der Punkte $h$ und $h^*$ und die völlig unbegründete Annahme gebracht, die drei geozentrischen Positionen lägen auf einem Großkreis. Er kann überhaupt nicht begreifen, wie es möglich sein soll, die Reduktion $c''-a''$ unmittelbar aus den Beobachtungen zu berechnen, da sie von zwei noch unbekannten Elementen abhängt, nämlich von den Abständen des Kometen und der Erde von der Sonne. Das ist allerdings richtig. Die Sache erklärt sich aber leicht daraus, dass von \emph{eben diesen} Elementen auch die Abweichung der mittleren Position von dem durch die äußeren Positionen gelegten Großkreis abhängt und gerade diese Abweichung dazu dient, $c''-a''$ zu bestimmen. Dass sie gewöhnlich sehr klein ist, tut hier nichts zur Sache, denn $c''-a''$ ist von derselben Ordnung. Die Darstellung, die \textsc{Güssmann} in § 7 von \textsc{Olbers}’ Gebrauch der orthografischen Projektion gibt, ist wieder ganz unrichtig. In seiner Fig. III haben die Punkte $b$, $i$ und $r$ mit unserer Sache gar nichts zu tun. Nur $z$ und $n$ haben damit zu tun. Die von $z$ nach $n$ gezogene Linie bezieht sich auf die erwähnte gedachte Position, die von $b$ nach $r$ gezogene dagegen nicht. Die Neigung der Projektion der Sichtlinie $z\,n$ bezeichnet \textsc{Olbers} mit $b''$, sodass dafür notwendig
$\operatorname{tang} b'' = \dfrac{\operatorname{tang} \gamma''}{\sin(A''-c'')}$
gilt. Dass nun $\dfrac{\operatorname{tang} \gamma''}{\sin(A''-c'')} = \dfrac{\operatorname{tang} \ObservedBeta''}{\sin(A''-a'')}$ gilt, folgt unmittelbar daraus, dass die Projektionen von $z\,n$ und $\mathfrak{S}\,m$ zusammenfallen. Wenn \textsc{Güssmann} danach in § 9 meint, \textsc{Olbers}’ Formel müsse eigentlich $\dfrac{t''\sin(a''-a')}{t'\sin(a'''-a'')}$ lauten, so ist das wiederum die Folge seiner unbegreiflichen Vermengung der Punkte $h$ und $h^*$. Wenn er aber den Unterschied zwischen \textsc{Olbers}’ Formel und diesem Ausdruck damit erklärt, dass \textsc{Olbers} alle drei Breiten $\ObservedBeta'$, $\ObservedBeta''$ und $\ObservedBeta'''$ berücksichtigt, die nicht alle drei zu einer Geraden passen können, und zwar an bestimmten Punkten dieser Geraden, von bestimmten Punkten der Erdsehne aus gesehen, dann urteilt \textsc{Güssmann} hier völlig richtig. Damit gibt er, ohne es zu wissen, den Grund dafür an, warum \textsc{Olbers}’ Formel richtig ist, während der Ausdruck $\dfrac{t''\sin(a''-a')}{t'\sin(a'''-a'')}$, der auf der Annahme einer geradlinigen Bewegung beruht, unrichtig ist. Denn tatsächlich lässt sich zeigen, dass diese letzte Formel nur dann eine brauchbare Näherung liefern kann, wenn die Abstände der Erde und des Kometen von der Sonne nicht sehr ungleich sind — doch hier ist nicht der Ort dafür. Dieser Umstand liegt zufällig in beiden Beispielen von \textsc{Olbers} vor. Nur deshalb ist der Unterschied zwischen den nach dem erwähnten Ausdruck berechneten Werten von $M$ und den nach \textsc{Olbers}’ Formel berechneten Werten nicht größer, als \textsc{Güssmann} in § 9 findet. Eben diesem Umstand ist zuzuschreiben, dass \textsc{Güssmann} in seinem Beispiel in § 24 nach jenem Ausdruck noch einigermaßen richtige Ergebnisse erhält. Die trivialen Bemerkungen, die \textsc{Güssmann} in § 11 macht, sind gewiss auch Dr. \textsc{Olbers} nicht entgangen. Dass $M$ bei kleinen Werten von $a'''-c''$ und $c''-a'$ ebenso sicher wie bei großen durch $\dfrac{t''\sin(c''-a')}{t'\sin(a'''-c'')}$ bestimmt wird, wenn man $c''$ mit vollkommener Genauigkeit kennt, ist allerdings wahr. Bei der begrenzten Genauigkeit aller auf Sinustafeln beruhenden Rechnungen bleibt jedoch beim berechneten Wert von $c''$ stets eine Unsicherheit. \textsc{Olbers} wollte
\SourceLeaf{494}
nur sagen, dass diese Unsicherheit, selbst wenn sie unter $0''1$ liegen sollte, in Fällen, in denen $a'''\!-c''$ und $c''\!-a'$ sehr klein sind, den Wert von $M'$ verfälschen könnte. Ebenso hat \textsc{Güssmann} mit seiner zweiten Bemerkung ganz recht, dass die Vernachlässigung des Unterschieds zwischen $c''$ und $a''$, den \textsc{Güssmann} mit $\delta$ bezeichnet, theoretisch selbst bei noch so kleinen Zwischenzeiten nicht zulässig ist. Unrecht hat er aber, wenn er glaubt, dass $\delta$ einen umso beträchtlicheren Teil von $a''\!-a'$ ausmacht, je kleiner die Zwischenzeit ist. Es verhält sich gerade umgekehrt, denn $a''\!-a'$ ist eine Größe erster, $\delta$ dagegen eine Größe zweiter Ordnung. Dass dessen ungeachtet auch bei sehr kleinen Zwischenzeiten die Vernachlässigung von $\delta$ einen bedeutenden Einfluss auf die Ergebnisse haben kann, wusste \textsc{Olbers} sehr wohl. Wenn er diese Vernachlässigung in diesem Fall dennoch zuließ, so geschah es, weil dieser Einfluss, abgesehen von seiner absoluten Größe, im Vergleich mit den Ungenauigkeiten aus den unvermeidlichen Beobachtungsfehlern umso kleiner wird, je kleiner die Zwischenzeiten sind. Wo die Unsicherheit der Beobachtungen ohnehin nur eine grobe Bestimmung der Bahn erlaubt, wäre es sehr verkehrt, sich durch ängstliche Rücksicht auf Umstände, deren Einfluss viel geringer ist als jene Unsicherheit, die Arbeit zu erschweren und damit eine Genauigkeit vorzutäuschen, die die Umstände nicht zulassen. Inwieweit es übrigens in solchen Fällen zweckmäßiger sein kann, die Rechnung lieber ganz zu unterlassen, als unsichere Ergebnisse herzuleiten, bleibt immer dem eigenen Urteil verständiger Astronomen anhand der jeweiligen Umstände überlassen.

Hiermit schließe ich meine Bemerkungen. Es genügt mir, die Angriffe besprochen zu haben, die \textsc{Güssmann} in theoretischer Hinsicht erhoben hat. Wer die Kürze und Bequemlichkeit von \textsc{Olbers}’ Methode mit den weitschweifigen Rechnungen vergleichen mag, die \textsc{Güssmann} anschließend anstellt und die doch nur grob angenäherte Ergebnisse liefern, wird sich auch darüber leicht ein Urteil bilden. Die Sache spricht für sich. Dass \textsc{Güssmann} in § 16 Herrn \textsc{von Pacassi} wegen der nach der Methode von \textsc{Boscovich} berechneten, aber ganz fehlerhaften Bahn des Kometen von 1779 damit rechtfertigt, dass es sehr leicht sei, in einer sehr \emph{weitläufigen} Rechnung einen Fehler zu begehen, der dann ein unrichtiges Ergebnis verursacht --- das ist doch ein beredter Beweis für den Vorzug, den die Methode von Dr. \textsc{Olbers} in dieser Hinsicht unbestreitbar hat. Es ist bekannt, dass sich viele, nicht nur einer, die Methode von \textsc{Olbers} mit Leichtigkeit zu eigen gemacht und damit Kometenbahnen sicher und richtig berechnet haben, darunter auch Menschen, die weder von Beruf Mathematiker waren noch beanspruchten, als besonders geübte Rechner zu gelten.
